%============================================================
\part{Theory}\label{part:theory}
%============================================================

\chapter{Introduction: Choosing the Time Structure of Cash as the Variable}

\section{The problem}

The standard variables for describing a firm are revenue, margin, and market share.
This text uses none of them. It describes a business model through the
\textbf{time structure of cash}.

There is a reason for the choice. What keeps a firm in existence is not a condition on
profit but a constraint on its cash balance, and that constraint is imposed path by path
rather than in expectation (Section~\ref{sec:objective}). Any framework that treats
survival and growth must take the time structure of cash as a first-class object.

The standard variables are not discarded.
\textbf{The income statement, the balance sheet, and the cash-flow statement appear as
projections of the object constructed here, each collapsing one of its indices}
(Table~\ref{tab:statements-projection}). This text does not posit quantities separate
from accounting; it keeps the indices that accounting collapses.

\section{Structure of the text, and how to read it}

The text has six parts.

\begin{center}
\begin{tabularx}{\textwidth}{@{}llX@{}}
\toprule
Part & Title & Content \\
\midrule
\ref{part:theory} & Theory & Definitions, propositions, proofs. The basis for everything that follows \\
\ref{part:catalog} & The space of maps & The whole space of $\Phi$ (28 types), the subspaces left under constraints, their limit, and the case where the information structure is supplied differently \\
\ref{part:method} & Method & The five stages that connect theory to observation, the treatment of selection bias, and a taxonomy of verification obstacles \\
\ref{part:results} & Empirics & How each theoretical quantity was measured, and tests in four domains \\
\ref{part:conclusion} & Conclusion & What was established and what remains open \\
\ref{app:proofs} & Appendices & Proofs of methodological claims, notation, correspondences with existing theory, data sources, pre-registration, revision history \\
\bottomrule
\end{tabularx}
\captionof{table}{The parts of this text.}
\label{tab:part-structure}
\end{center}

Parts~\ref{part:theory} and~\ref{part:catalog} are deductive;
Part~\ref{part:method} supplies the procedure that connects theory to evidence.

Part~\ref{part:catalog} proceeds in three stages.
It first enumerates the space generated by the degrees of freedom of $\Phi$
(Chapter~\ref{ch:catalog}), then derives the subspace that survives constraints on capital
and labour (Chapter~\ref{ch:solo}), and finally treats the limit in which the operator is
removed (Chapter~\ref{part:unmanned}). All three are deduction, not evidence.

The absence of a separate part on applications is deliberate.
An application is warranted once the theory has passed through evidence, and the empirical
work here is not complete. Deduction under constraints is not application; it is
\textbf{the operation of cutting the space of maps with a constraint}, and it therefore
belongs in Part~\ref{part:catalog}.

Chapter~\ref{ch:instantiation} opens Part~\ref{part:results} and tabulates, for every symbol
introduced in the theory, the value that could be measured and the fact that could not.
Each later chapter also opens with a table of the theoretical quantities it handles, so that
one can trace where each symbol was measured.

Depending on the reader's interest, the following routes are suggested.

\begin{itemize}
\item \textbf{Theory alone}: Parts~\ref{part:theory} and~\ref{part:catalog}.
Proofs of the propositions are in the body; only the proofs of methodological claims are
collected in Appendix~\ref{app:proofs}
\item \textbf{Practical design}: Chapter~\ref{sec:growth} and Part~\ref{part:catalog}
\item \textbf{The magnitudes the theoretical quantities actually take}: Chapter~\ref{ch:instantiation}
\item \textbf{Survey method}: Parts~\ref{part:method} and~\ref{part:results}
\end{itemize}

Numerical examples accompany most sections. The notation is listed in
Appendix~\ref{app:notation}.

\section{What this text does not claim}

To forestall misreading, the scope is limited in advance.

First, this is \textbf{not a predictive theory}. Most of the implications fixed in advance
and tested in Part~\ref{part:results} were rejected. The framework works as a descriptive
instrument, but this text has no established result about what the quantities it describes
predict.

\textbf{Twenty-two implications were nevertheless fixed in advance.} There are two reasons.

The first is that a descriptive framework \textbf{can supply an after-the-fact account of
any fact whatever}. That something can be explained is no evidence that the framework is
right. Unless the implications are fixed in writing before the observations and the ones
that miss are left standing as misses, there is no way to distinguish the framework working
from the author having reasoned backwards. The places where an after-the-fact story was
nearly adopted are listed in Section~\ref{sec:prereg-worked}.
\textbf{The fixed implications are not claims of the framework; they are a constraint
imposed on the author's own interpretation.}

The second is that \textbf{an implication failing and the framework being rejected are
different events}. Whether the descriptions $\CCC$, $\kappa$ and $\Phi$ work is one
question; what those quantities predict is another. A failed implication refutes a
conjecture about the latter, not the former (Section~\ref{sec:two-rejections}).

Second, this text is \textbf{not normative}. It does not argue which $\Phi$ is desirable,
only which $\Phi$ is feasible under the constraints.

Third, \textbf{most of the individual components have counterparts in existing theory}.
Appendix~\ref{app:priorwork} is divided into three layers separating the existing theory
that explains each term, the part that connects them, and the organization contributed here.
The contribution of this text lies in the connection, not in the novelty of the elements.

%============================================================
\chapter{The Business Model: \texorpdfstring{$\Phi$}{Phi} and the Generation of Cash Flow}
\label{sec:phi}
%============================================================

\section{Primitive sets}

\begin{definition}[Primitive sets]
Four sets are taken as primitive.
\begin{align}
N &= \{0,1,\dots,n\} && \text{the set of parties; } 0 \text{ is the firm in question} \\
T &= [0,H] \subset \R_{\geq 0} && \text{time} \\
(\Omega,\mathcal{F},\{\mathcal{F}_t\}_{t\in T},\Prob) && &\text{state space and filtration} \\
\mathcal{D} && &\text{the space of deliveries (goods and services)}
\end{align}
The parties are grouped into customers $C$, suppliers $S$, capital providers $K$,
labour $L$, and third parties $R$ who assume risk. These are subsets of $N$ and cover
$N\setminus\{0\}$. An individual party is written as an element of $N$, as in $i\in C$.
When a group symbol is used directly as an index it denotes the sum over the parties in
that group; for instance $\kappa_C=\sum_{i\in C}\kappa_i$.
$R$ is where degree of freedom (3) moves the paying party, which is the transfer operation
of Section~\ref{sec:three-ops}.

Regulation and commercial custom are not counterparties to cash and are therefore not in
$N$. They constrain the range from which $\Phi$ may be chosen and are treated as the
institutional layer in Section~\ref{sec:layers}.
\end{definition}


\section{Schedules and spaces}
\label{sec:schedules}

The relationship with party $i$ is described as a pair of schedules.
\begin{align}
\delta_i : T\times\Omega \to \mathcal{D} && \text{delivery (value handed from the firm to $i$)} \\
\pi_i : T\times\Omega \to \R && \text{settlement (payment from $i$ to the firm)}
\end{align}
A valuation map $v:\mathcal{D}\to\R$ measures deliveries in monetary units.

The collections of such tuples are taken as spaces.
\begin{equation}
\Delta = \bigl\{\, (\delta_i)_{i\in N} \ \bigm| \ \delta_i : T\times\Omega\to\mathcal{D} \,\bigr\},
\qquad
\Pi = \bigl\{\, (\pi_i)_{i\in N} \ \bigm| \ \pi_i : T\times\Omega\to\R \,\bigr\}
\label{eq:spaces}
\end{equation}
$\Delta$ is called the \emph{space of delivery schedules} and $\Pi$ the
\emph{space of settlement schedules}.

\section{The map from delivery to settlement}

\begin{definition}[Business model]
Using $\Delta$ and $\Pi$ from~\eqref{eq:spaces}, a business model is defined as the map
\begin{equation}
\Phi : \Delta \longrightarrow \Pi
\label{eq:phi}
\end{equation}
\end{definition}

\begin{remark}[$\Phi$ and $\phi$ are different objects]
\label{rem:Phi-vs-phi}
Capital $\Phi$ denotes the map~\eqref{eq:phi} itself, that is, the \textbf{structure}.
Lower-case $\phi$ denotes the \textbf{size of the surplus} that arises in one period under
that structure; it is decomposed into three sources in~\eqref{eq:surplus} of
Chapter~\ref{sec:surplus}. $\Phi$ fixes the shape of the correspondence between delivery
and settlement; $\phi$ measures the share the firm retains under that shape. No general
relation between the two is given here. How the choice of $\Phi$ moves each term of $\phi$
is discussed term by term in Chapter~\ref{sec:surplus}.
\end{remark}


\section{Generating the cash flow}
\label{sec:cashflow}

Once $\Phi$ in~\eqref{eq:phi} is fixed, the settlement with each party is fixed. This is
called the cash flow.

\begin{definition}[Cash flow]
Given a delivery $\delta\in\Delta$, the map $\Phi$ returns as its image the tuple of
settlements $\pi = \Phi(\delta) \in \Pi$. An element of $\Pi$ is a tuple of functions
indexed by party, so its $i$-th component $\pi_i$ is a function on $T\times\Omega$. Write
\begin{equation}
X(i,t,\omega) \ \equiv\ \pi_i(t,\omega), \qquad \pi = \Phi(\delta)
\label{eq:X}
\end{equation}
Then $X(i,t,\omega)$ is the net movement of cash between the firm and party $i$ at time $t$
in state $\omega$. Inflows to the firm are positive. Payments to suppliers or to labour
appear as $\pi_i<0$.
\end{definition}

\textbf{$\Phi$ takes exactly one argument, the tuple of deliveries $\delta$; it does not
take $t$ or $\omega$.} The dependence on time and state is carried by the $\pi_i$ that
$\Phi$ returns.

\textbf{$X$ is the image of $\Phi$.} Replacing $\Phi$ produces a different $X$ from the
same delivery $\delta$. This is the central operation of the text.

Equation~\eqref{eq:X} carries three indices: who, when, and in which state. The usual
financial statements are projections that collapse this object along one index.

\begin{center}
\begin{tabularx}{\textwidth}{@{}llX@{}}
\toprule
Statement & Operation & Content \\
\midrule
Income statement & aggregation over $t$ & includes accrual adjustments \\
Balance sheet & a cross-section at fixed $t$ & unsettled balances \\
Cash-flow statement & differences along $t$ & the actual movement of cash \\
\bottomrule
\end{tabularx}
\captionof{table}{The financial statements are projections that collapse the three indices of~\eqref{eq:X}.}
\label{tab:statements-projection}
\end{center}

\begin{remark}[Every quantity here is a level]
\label{rem:level-vs-change}
Every quantity introduced after~\eqref{eq:X} --- including $\kappa_i$, $\CCC$ and $\phi$ ---
is a level. Each is defined as a value at a time $t$, and no notation for change is
provided.

This is not a presentational choice but \textbf{a limitation}. As Part~\ref{part:results}
confirms repeatedly, the levels of these quantities differ by two orders of magnitude across
industries and product forms, and \textbf{cross-sectional comparison cannot be interpreted}.
What can be read is change over time within one object.

The text has no theoretical apparatus for change and deals with it case by case in the
empirical work (Section~\ref{sec:level-change-gap}).
\end{remark}

\begin{example}[The three indices made concrete]
\label{ex:triple-index}
Consider a service priced at 1{,}000 yen a month and supplied to 100 customers.
If customer $i\in C$ pays at the end of each month, then
\[
X(i,\ t,\ \omega) = \begin{cases} 1{,}000 & t \in \{1,2,\dots\}\ \text{(month end)} \\ 0 & \text{otherwise}\end{cases}
\]
which does not depend on $\omega$. If cancellation is possible, $\omega$ contains the
cancellation time and $X$ does depend on $\omega$.

With a single supplier, $S=\{s\}$, and a server cost of 30{,}000 yen a month,
$X(s,t,\omega)=-30{,}000$. Written with the group index,
$\sum_{i\in S}X(i,t,\omega)=-30{,}000$.
\end{example}


\section{Degrees of freedom of \texorpdfstring{$\Phi$}{Phi}}
\label{sec:dof}

The features that specify $\Phi$ in~\eqref{eq:phi} are organized into six degrees of
freedom. They divide into those that correspond to the three indices of~\eqref{eq:X} and
those that do not.

\paragraph{Those that move an index.}
\begin{enumerate}[label=(\arabic*)]
\item \textbf{Timing}: the order of $\operatorname{supp}(\pi)$ and $\operatorname{supp}(\delta)$. Moves $t$. Fixes the sign of $\kappa$.
\setcounter{enumi}{2}
\item \textbf{Paying party}: the beneficiary, a third party, or the opposite side of the market. Moves the index $i$.
\item \textbf{State dependence}: whether $\pi$ is measurable with respect to $\omega$. Moves $\omega$, and fixes where risk resides.
This is exactly the $\mathcal{G}$-measurability of Chapter~\ref{sec:info}; the range the
designer of $\Phi$ may choose from is bounded by $\mathcal{G}$
(Proposition~\ref{prop:implementable}, Remark~\ref{rem:G-limits-dof4}).
\end{enumerate}

\paragraph{That which fixes the functional form of the map.}
\begin{enumerate}[label=(\arabic*)]
\setcounter{enumi}{1}
\item \textbf{Form of dependence}: how $\pi$ depends on $\delta$.
The linear forms collapse into a single family.
\begin{equation}
\pi = F + p\,\delta
\label{eq:two-part}
\end{equation}
$F$ is the part independent of the quantity delivered and $p$ is the unit price.
$p=0$ gives a flat fee, $F=0$ gives usage pricing, and both non-zero gives a two-part
tariff. A non-linear form is outcome-contingent pricing $\pi = h(y)$, where $y$ is the
outcome and $h$ turns the outcome into a settlement; whether $y$ is observable belongs to
degree of freedom (4) (Chapter~\ref{sec:info}).
\end{enumerate}

\begin{remark}[The typeface of $F$, and the use of $c$]
\label{rem:dof2-symbols}
The $F$ in~\eqref{eq:two-part} is the fixed part of a tariff and is a different quantity
from the $\sigma$-algebra $\mathcal{F}$ on the state space. The two are distinguished by
typeface.

Throughout this text $c$ denotes marginal cost and is always written with its argument,
$c(\delta)$ (Definition~\ref{def:cost}). A flat fee is written $F$.
\end{remark}

\paragraph{Properties on the side of the domain \texorpdfstring{$\Delta$}{Delta}.}
\begin{enumerate}[label=(\arabic*)]
\setcounter{enumi}{4}
\item \textbf{Separation of right from exercise}: whether the contractual right $\dbar$ and realized use $\delta$ can diverge.
\item \textbf{Repetition}: one-off, auto-renewing, or with a committed term. Fixes the transaction frequency $\nu$.
\end{enumerate}

\begin{remark}[Three groups]
\label{rem:dof-groups}
(1), (3) and (4) move the indices of~\eqref{eq:X}: they \textbf{change where $X$ is placed
for the same $\delta$}. (2) fixes the form of the map itself
(Remark~\ref{rem:dof2-symbols}), and (5) and (6) are properties of the domain $\Delta$.

The three risk operations of Section~\ref{sec:three-ops} are the operations that move
(1), (3) and (4).
\end{remark}

Of these, (1) and (2) respectively fix the sign of $\kappa$ and the room for cognitive
surplus. These two axes give a coarse partition of the space of $\Phi$
(Figure~\ref{fig:quadrant}). \textbf{The two are taken from different groups}: (1) moves an
index and (2) fixes a functional form.

\textbf{The two axes of the partition and the degrees of freedom used to assign a family
are different things.} The assignment order of Section~\ref{sec:type-assignment} uses (3),
(4) and (6); (1) and (2) do not appear. The former are axes for viewing the space of $\Phi$
coarsely, the latter a rule for assigning an observed $\Phi$ uniquely; the purposes differ.

\begin{figure}[htbp]
\centering
\insertfig{kappa-schedule}
\caption{A two-axis partition of $\Phi$. The vertical axis fixes the sign of $\kappa$, the horizontal axis the room for cognitive surplus.}
\label{fig:quadrant}
\end{figure}

\section{Three risk operations}
\label{sec:three-ops}

Because there are three indices, there are only three operations on $X$.
The correspondence is an organization contributed here.

\begin{center}
\begin{tabular}{@{}clll@{}}
\toprule
Index moved & Name & Invariant & Examples \\
\midrule
$t$ & timing & total & advances, instalments, leases \\
$i$ & transfer & total and variance & factoring, buying insurance \\
$\omega$ & pooling & total & underwriting, diversified investment \\
\bottomrule
\end{tabular}
\captionof{table}{Three risk operations: the index moved, the invariant, and typical examples.}
\label{tab:three-ops}
\end{center}

As Remark~\ref{rem:dof-groups} states, these three are the operations that move degrees of
freedom (1), (3) and (4): timing moves $t$, transfer moves $i$, and pooling moves $\omega$.

Of the three, only pooling actually reduces variance. This is a consequence of probability
theory and holds independently of the structure of $\Phi$.
\textbf{It is quoted here, not derived.}

\begin{proposition}[Variance reduction by pooling]
\label{prop:pooling}
For identically distributed risks $X_1,\dots,X_n$ with variance $\sigma^2$ and pairwise
correlation $\rho$, \textbf{provided $\rho$ does not depend on $n$},
\begin{equation}
\operatorname{Var}\!\left[\frac{1}{n}\sum_{j=1}^{n} X_j\right]
= \frac{\sigma^2}{n}\bigl(1+(n-1)\rho\bigr)
\xrightarrow[n\to\infty]{} \rho\,\sigma^2 .
\label{eq:pooling}
\end{equation}
\end{proposition}

\begin{proof}
By bilinearity of covariance,
\[
\operatorname{Var}\Bigl[\sum_j X_j\Bigr]
= \sum_j \operatorname{Var}[X_j] + \sum_{j\neq k}\operatorname{Cov}[X_j,X_k]
= n\sigma^2 + n(n-1)\rho\sigma^2 .
\]
Dividing both sides by $n^2$ gives the first equality. As $n\to\infty$,
$\frac{\sigma^2}{n}\to 0$ and $\frac{(n-1)\rho\sigma^2}{n}\to\rho\sigma^2$, which gives the
limit.
\end{proof}

\begin{remark}[Independence is a finite resource]
\label{rem:independence-finite}
The limit in~\eqref{eq:pooling} is $\rho\sigma^2$, which is non-zero whenever $\rho>0$.
Pooling consumes independence, and independence is a finite resource.
Underwriting limits on earthquake insurance, the sensitivity of platforms to the business
cycle, and systemic risk among financial institutions all come from the same term of this
equation.
\end{remark}

\begin{remark}[Endogeneity of $\rho$]
\label{rem:rho-endogenous}
Proposition~\ref{prop:pooling} treats $\rho$ as an exogenous constant. The assumption is
hard to sustain.

The mechanism that makes $\rho$ endogenous is not the exhaustion of poolable risks but the
fact that \textbf{the act of pooling itself creates correlation}. The example of
\cite{wagner2010} is clear: if two parties each take half of two risks and diversify
perfectly, each party's individual probability of failure falls, but
\textbf{their portfolios become perfectly correlated}. The nature of the risk is unchanged,
yet correlation between the parties appears.

Hence $\rho$ is a function of the number of poolers and the overlap of their holdings, and
the limit $\rho\sigma^2$ in~\eqref{eq:pooling} can rise as pooling proceeds. The proof is
formally correct, but \textbf{its range of application is confined to cases with few poolers
and little overlap}. This dependence is not treated here
(Section~\ref{sec:independence-assumption}).
\end{remark}

\begin{remark}[Social optimality of pooling]
\label{rem:pooling-social}
The mechanism of Remark~\ref{rem:rho-endogenous} implies that the individual and the
collective optimum can diverge. \cite{ibragimov2011} derives conditions under which the
diversification optimal for an intermediary is not socially optimal and restricting risk
sharing is preferable.

This text is written from the standpoint of an individual operator and does not treat the
effect of the choice of $\Phi$ on the system. Family 5-4 (underwriting) and family 6-2
(escrow) in Part~\ref{part:catalog} make pooling their business and are the types where
this point bears directly.
\end{remark}

\begin{example}[How fast correlation bites]
With $\sigma=100$ and $\rho=0$, at $n=100$ the standard deviation falls to
$100/\sqrt{100}=10$. With $\rho=0.3$, however,
\[
\operatorname{Var}=\frac{100^2}{100}\bigl(1+99\times 0.3\bigr)=305.5,\qquad
\text{standard deviation}\approx 17.5 .
\]
Even as $n\to\infty$ the standard deviation does not fall below
$\sqrt{0.3}\times 100 \approx 54.8$.
At a correlation of $0.3$ --- a modest value --- the benefit of pooling is more than halved.
\end{example}

\subsection{Cost fixed by technology}
\label{sec:production}

Producing a delivery $\delta$ costs something. Because this text does not treat the
composition of inputs, \textbf{production technology appears only as a correspondence
$C_{\mathrm{prod}}$ from delivery to cost}.

\textbf{$C_{\mathrm{prod}}$ is given independently of $\Phi$.} Different $\Phi$ can be
placed on the same technology, and the same $\Phi$ on different technologies. What this
text varies is $\Phi$, not the technology.

\begin{definition}[Splitting the cost]
\label{def:cost}
Total cost is split in two.
\begin{equation}
C(\delta,\Phi) = \underbrace{C_{\mathrm{prod}}(\delta)}_{\text{fixed by technology}}
 + \underbrace{C_{\mathrm{admin}}(\Phi,\delta)}_{\text{created by the contract}}
\label{eq:cost-split}
\end{equation}
For the first term, write
\begin{align}
C_{\mathrm{prod}}(0) && &\text{fixed cost} \\
c(\delta) = \frac{dC_{\mathrm{prod}}}{d\delta} && &\text{marginal cost}
\end{align}
\end{definition}

\begin{remark}[No shape is assumed for $c(\delta)$]
Neither monotonicity nor convexity is assumed for $c(\delta)$. Cost structures differ
greatly across industries, and imposing one shape would narrow the range of application.

The price is that \textbf{the optimal quantity delivered $\delta$ cannot be fixed at an
interior point}. The capacity constraint~\eqref{eq:cap-simplified} supplies an upper bound,
but the point short of it at which cost balances price does not follow from the assumptions
made here. Where the level of $\delta$ matters later --- the allocation in
Chapter~\ref{ch:solo}, for instance --- it is treated as a corner of the constraint set.
\end{remark}

The split is meaningful because the second term is real.
$C_{\mathrm{admin}}$ arises from all six degrees of freedom of $\Phi$
(Section~\ref{sec:dof}).

\begin{center}
\begin{tabularx}{\textwidth}{@{}llX@{}}
\toprule
Kind & Route & Content \\
\midrule
Setup & conclusion & the cost of agreeing $\Phi$ with the counterparty; small for standard contracts, large for negotiated ones \\
Operation & measurement & under $\pi=p\,\delta$, $\delta$ must be metered and billed \\
 & rights management & administering $\dbar$: memberships, tracking balances \\
 & ongoing handling & auto-renewal, cancellation \\
Securing performance & collection and bad debt & under $\tau_C>0$, invoicing, chasing, and losses on non-collection \\
 & verification & under $\pi=h(y)$, confirming $y$ \\
 & dispute & disputes arise even when the contract is written \\
Outsourcing & intermediation & settlement, delivery, customer acquisition performed by others \\
\bottomrule
\end{tabularx}
\captionof{table}{Routes through which $C_{\mathrm{admin}}$ arises.}
\label{tab:cadmin-routes}
\end{center}

\begin{remark}[Advances cost something too]
$\tau_C<0$ (an advance) requires no collection cost, but it does create custody of the
funds held, refund handling, and regulatory compliance. The deposit obligation under the
Payment Services Act, treated in Section~\ref{sec:bootstrap}, is an example.
\textbf{$C_{\mathrm{admin}}>0$ whatever the sign of $\tau$.}
\end{remark}

\begin{remark}[Conditions under which $C_{\mathrm{admin}}$ is ignored]
\label{rem:cadmin-negligible}
From here on $C_{\mathrm{admin}}$ is not written explicitly. Two conditions make it
negligible.

\begin{enumerate}[label=(\arabic*)]
\item \textbf{Operations are systematized.}
Once the machinery is built, the marginal cost per transaction of measurement, rights
management and ongoing handling approaches zero.
\item \textbf{No intermediary is used.}
Intermediation fees are proportional to transaction value and do not vanish with
systematization. Chapter~\ref{sec:platform-fee} measures this cost and maps it onto the
kinds of $C_{\mathrm{admin}}$ in Remark~\ref{rem:ladder-is-cadmin}.
\end{enumerate}

The range in which the conditions hold is limited. They hold for the solo business of
Chapter~\ref{ch:solo} and for family 2 in Part~\ref{part:results}, but
\textbf{not for the small firms treated in Chapter~\ref{part:industry}}.
Where invoicing and chasing are done by hand, $C_{\mathrm{admin}}$ cannot be ignored.

The price of the omission should be stated. Because $C_{\mathrm{admin}}$ is dropped, this
text \textbf{cannot treat the route by which the choice of $\Phi$ feeds back into $\phi$
through cost}. Equation~\eqref{eq:cost-split} is not used again; only $C_{\mathrm{prod}}$
appears. A relation such as ``lengthening $\tau_i$ raises collection cost'' cannot be
written within this framework.
\end{remark}

What is called a business model here is therefore the freedom to choose different
contractual forms on an identical cost structure. That freedom is not complete: the choice
of $\Phi$ feeds back into cost through $C_{\mathrm{admin}}$.
Where $C_{\mathrm{prod}}$ appears is collected in Remark~\ref{rem:cost-usage}.

\begin{remark}[Switching cost and the value of waiting]
\label{rem:switching-cost}
Treating $\Phi$ as a choice variable amounts to \textbf{setting the switching cost to zero}.
The cost of changing from one $\Phi$ to another belongs to the setup term of
$C_{\mathrm{admin}}$ and is not generally zero.

In the theory of irreversible investment (\cite{dixit1994}), a decision is irreversible not
because it cannot physically be undone but because \textbf{undoing it is expensive and locks
in a particular path}. Writing $C_{\mathrm{sw}}$ for the switching cost and $\Delta\phi$ for
the gain in surplus from changing $\Phi$:

\begin{center}
\begin{tabularx}{\textwidth}{@{}lX@{}}
\toprule
$C_{\mathrm{sw}} = \infty$ & no change is possible; the first choice is final \\
$C_{\mathrm{sw}} \ll \Delta\phi$ & the assumption made here: $\Phi$ is freely chosen \\
$C_{\mathrm{sw}} \sim \Delta\phi$ & there is value in delaying the choice \\
\bottomrule
\end{tabularx}
\captionof{table}{The size of the switching cost $C_{\mathrm{sw}}$ and the freedom to choose $\Phi$.}
\label{tab:switching-cost-cases}
\end{center}
In the third case \textbf{waiting has value}: it can be optimal not to fix $\Phi$ until the
uncertainty resolves. This implication is not treated here. For the solo business of
Chapter~\ref{ch:solo}, where the switching cost is relatively large, the point may not be
negligible.
\end{remark}

\begin{remark}[Where $C$ appears]
\label{rem:cost-usage}
The quantities of Definition~\ref{def:cost} are used repeatedly.
$\phi_{\mathrm{prod}}$ is defined as an advantage in $c(\delta)$ (Chapter~\ref{sec:surplus});
a near-zero marginal cost is the condition that relaxes the capacity constraint
(Chapter~\ref{ch:solo}); and the fact that both $C$ and $\Phi$ can be specified in advance
is the condition for replicability (Chapter~\ref{part:unmanned}). The physical layer refers
to these properties of $C$ (Section~\ref{sec:layers}).
\end{remark}

\begin{example}[Identical $C_{\mathrm{prod}}$, different $\Phi$]
The same business software can be supplied in three ways.
\begin{center}
\begin{tabular}{@{}llr@{}}
\toprule
$\Phi$ & Settlement & $\kappa_C$ (mid-period, one customer) \\
\midrule
Perpetual licence & 240{,}000 yen on installation & \textminus 120{,}000 yen \\
Annual & 60{,}000 yen at the start of each year & \textminus 30{,}000 yen \\
Monthly in arrears & 5{,}000 yen at each month end & +2{,}500 yen \\
\bottomrule
\end{tabular}
\captionof{table}{Three choices of $\Phi$ on an identical $C_{\mathrm{prod}}$, and the resulting $\kappa_C$.}
\label{tab:same-f-different-phi}
\end{center}
The technology is identical and so is the marginal cost. Only $\Phi$ differs, yet the sign
of $\kappa$ reverses and the working capital required changes.
\end{example}

\section{Divergence between right and exercise}
\label{sec:breakage}

Degrees of freedom (2) and (5) are linked. When $\pi$ is proportional to $\delta$ the
customer pays for what is used, so right and exercise coincide. Only when $\pi$ is
independent of $\delta$ can the two diverge.

When degree of freedom (5) of~\eqref{eq:phi} is present, define the following.

\begin{definition}[Right and exercise]
\label{def:breakage-def}
Let $\dbar$ be the ceiling the contract grants the customer and $\delta$ the use actually
realized. Write the divergence as $\dbar-\delta \geq 0$ and define the gross margin arising
from it as
\begin{equation}
\phi_{\mathrm{cog}} = p\cdot\E[\dbar-\delta]
\label{eq:cog}
\end{equation}
\end{definition}

\begin{remark}[Which contractual forms admit divergence]
\label{prop:breakage}
Equation~\eqref{eq:cog} can be positive only when the contract states $\dbar$ explicitly.

Under usage pricing $\pi(t)=p\,\delta(t)$, the customer's total payment is $p\int\delta$.
No use means no payment, so the contract need not state a ceiling $\dbar$ and the difference
$\dbar-\delta$ is undefined.

Under a flat fee, where $\pi$ does not depend on $\delta$, the consideration is fixed
irrespective of use, so the contract states the extent of the right $\dbar$. Here
$\delta<\dbar$ can be realized. Under a two-part tariff $\pi=F+p\,\delta$, divergence
arises only over the range corresponding to the fixed part $F$.

\textbf{This is a consequence of Definition~\ref{def:breakage-def}, not an independent
claim.} It presumes that $\dbar$ is given exogenously. When the customer chooses $\dbar$
from a menu, the two are determined jointly and the meaning of the difference changes
(Section~\ref{sec:independence-assumption}).
\end{remark}

\begin{remark}[The general form of the capacity constraint, and the simplification used here]
\label{rem:cap-simplification}
Capacity on the supply side covers the occupancy of a facility, the number of machines, the
disposable time of staff, and so on. Since $\delta_i$ is a function on $T\times\Omega$ by
the definitions of Chapter~\ref{sec:phi}, the general form of the capacity constraint is
evaluated at each instant.
\begin{equation}
\sum_i \delta_i(t,\omega) \ \leq\ \mathrm{Cap}(t,\omega)
\qquad \forall t\in T,\ \forall\omega\in\Omega
\label{eq:cap-general}
\end{equation}
Capacity itself depends on time because opening hours, equipment run-times and staffing all
vary over time.

\textbf{From here on $\mathrm{Cap}$ is treated as a constant}, that is,
\eqref{eq:cap-general} is used in the form integrated over the whole period,
\begin{equation}
\sum_i \delta_i \ \leq\ \mathrm{Cap}
\label{eq:cap-simplified}
\end{equation}
where $\delta_j$ denotes a total.

The simplification is justified when demand is distributed over time in roughly the same way
across members and capacity does not vary by time of day. When neither holds, only the peak
instant binds and~\eqref{eq:cap-simplified} is stronger than necessary.
Chapter~\ref{part:industry} treats an observation with which this simplification does not
agree.
\end{remark}

\begin{proposition}[Membership under a capacity constraint]
\label{prop:capacity}
Suppose the capacity constraint~\eqref{eq:cap-simplified} is present and
\textbf{the number of members $n$ and average use $\E[\delta]$ are determined independently}
(the number of members is written $n$ because $m$ is reserved for the margin).
Then a flat membership $\pi_j=F$ is feasible only if
\begin{equation}
\sum_{j}\E[\delta_j] \ \ll\ \sum_{j}\dbar_j
\end{equation}
that is, low utilization is a condition of feasibility.
\end{proposition}

\begin{proof}
Let $n$ be the number of members and $\dbar$ each member's contractual ceiling. If every
member exercised the right, demand would be $n\dbar$. Feasibility requires
$n\,\E[\delta] \leq \mathrm{Cap}$. On the other hand, since members join because they value
$\dbar$, normally $n\dbar > \mathrm{Cap}$ (otherwise the capacity constraint does not bind).
Combining the two gives $\E[\delta] \leq \mathrm{Cap}/n < \dbar$.
\end{proof}

\begin{example}[Two kinds of membership]
A facility with room for 100 has 1{,}000 members and charges 8{,}000 yen a month.
If $\dbar$ is ``use it every day'', $n\dbar$ far exceeds capacity. At an actual utilization
of 10\%, $n\E[\delta]=100$ and the constraint is just satisfied.
\textbf{If utilization rises to 30\%, the business fails physically.}

By contrast, where the fee grants a right to purchase discounts there is no capacity
constraint, and higher utilization is favourable because it enlarges the purchasing volume.
Outwardly both belong to family 2-1, but the shape of the feasible region differs.
\end{example}

%============================================================
\chapter{The Credit Position: \texorpdfstring{$\kappa$}{kappa} and the Lag \texorpdfstring{$\tau$}{tau}}
\label{sec:kappa}
%============================================================

\section{The gap between delivery and settlement}

Write the cumulants of the delivery $\delta_i$ of Section~\ref{sec:schedules} and of the
settlement $\pi_i$ fixed by $\Phi$ in~\eqref{eq:phi} (with $\pi=\Phi(\delta)$) as
\begin{equation}
D_i(t) = \int_0^t v(\delta_i(s))\,ds, \qquad
P_i(t) = \int_0^t \pi_i(s)\,ds
\label{eq:cum}
\end{equation}
$D_i$ is value given and $P_i$ is consideration received.

\begin{definition}[Credit position]
The credit position towards party $i$ is defined as
\begin{equation}
\kappa_i(t) = D_i(t) - P_i(t)
\label{eq:kappa}
\end{equation}
By~\eqref{eq:cum},
\[
\kappa_i(t) = \int_0^t v(\delta_i(s))\,ds - \int_0^t \pi_i(s)\,ds
\]
so \textbf{$\kappa_i$ is a function of the $\Phi$ of~\eqref{eq:phi}}: replacing $\Phi$ for
the same delivery $\delta$ changes $\kappa_i$. A positive $\kappa_i$ means the firm is
extending credit to $i$ (delivered, not yet collected); a negative $\kappa_i$ means $i$ is
extending credit to the firm (received, not yet delivered).
\end{definition}

The sum of credit positions plus inventory is used repeatedly below.

\begin{definition}[Working capital]
\label{def:wc}
With inventory $\mathrm{Inv}(t)\geq 0$, define
\begin{equation}
W(t) = \sum_{i\in N}\kappa_i(t) + \mathrm{Inv}(t)
\label{eq:wc}
\end{equation}
Since $\kappa_i$ is a function of $\Phi$, so is $W$.
\end{definition}

The $W$ of~\eqref{eq:wc} is the amount of money tied up in the business; dividing it by the
rate of sales gives $\CCC$ in Chapter~\ref{sec:growth}.

The merit of~\eqref{eq:kappa} is that it unifies.

\begin{center}
\begin{tabular}{@{}llll@{}}
\toprule
Account & Party & Sign & Meaning \\
\midrule
Receivables & $C$ & $\kappa_C>0$ & the firm extends credit to the customer \\
Advances received, contract liabilities & $C$ & $\kappa_C<0$ & the customer extends credit to the firm \\
Payables & $S$ & $\kappa_S<0$ & the supplier extends credit to the firm \\
Advances paid & $S$ & $\kappa_S>0$ & the firm extends credit to the supplier \\
\bottomrule
\end{tabular}
\captionof{table}{Four accounts unified by~\eqref{eq:kappa}.}
\label{tab:kappa-accounts}
\end{center}

These are not four separate accounts but one quantity with different signs and indices.

\subsection{Existing theory for each term}
\label{sec:kappa-literature}

Each term of $\kappa_i$ has an established theory behind it. This text takes them as its
foundation.

\begin{center}
\begin{tabularx}{\textwidth}{@{}llX@{}}
\toprule
Party and sign & Phenomenon & Explanatory theory \\
\midrule
$\kappa_C>0$ (inter-firm) & selling on account & four motives for trade credit: financing advantage, price discrimination, quality assurance, liquidation advantage \cite{petersen1997,schwartz1974} \\
$\kappa_C<0$ (inter-firm) & advance payment & reverse trade credit: supplier finance, transaction assurance, bargaining power \cite{mateut2014,daripa2011} \\
$\kappa_C<0$ (to consumers) & prepaid contracts & mental accounting, evidence on breakage \cite{nber2026pnbl} \\
$\kappa_S<0$ & payables & the mirror image of trade credit \\
$\kappa_L<0$ & wages in arrears & the bonding argument: discipline through deferred pay \cite{akerlof1986} \\
$\kappa_K<0$ & borrowing & standard corporate finance \\
$\kappa_K^{i}$ & equity (residual claim) & party-dependent; Definition~\ref{def:kappa-K} \\
\bottomrule
\end{tabularx}
\captionof{table}{Existing theory explaining each term of $\kappa_i$.}
\label{tab:kappa-literature}
\end{center}

\begin{remark}[The four literatures do not cite one another]
\label{rem:disjoint-literatures}
The theories in Table~\ref{tab:kappa-literature} developed in different fields to answer
different questions.

\begin{center}
\begin{tabularx}{\textwidth}{@{}llX@{}}
\toprule
Theory & Field & Central question \\
\midrule
Trade credit & corporate finance & why a supplier can lend more cheaply than a bank \\
Bonding argument & labour economics & why involuntary unemployment exists \\
Signalling & economics of information & how ability is disclosed \\
Work on prepaid contracts & consumer behaviour & why consumers pay first \\
\bottomrule
\end{tabularx}
\captionof{table}{The fields and central questions of the theories explaining each term of $\kappa_i$.}
\label{tab:disjoint-theories-questions}
\end{center}

Because the questions differ, these literatures do not cite one another. Yet
in~\eqref{eq:cash} \textbf{the customer's prepayment, the supplier's trade credit, the
worker's deferred wage, and the investor's equity all appear as terms of the same sum}.

Each theory explains one term of the sum; none has a framework that treats the sum. The
contribution here is that integration, not the explanation of any individual term.
\end{remark}

The relation between the cumulants is shown in Figure~\ref{fig:kappa}.

\begin{figure}[htbp]
\centering
\insertfig{phi-quadrant}
\caption{The credit position measured against cumulative delivery $D(t)$. Where the settlement curve $P(t)$ lies above $D(t)$, $\kappa<0$; where it lies below, $\kappa>0$.}
\label{fig:kappa}
\end{figure}

\subsection{Both sides of a transaction: $\kappa$ is antisymmetric}
\label{sec:kappa-antisym}

Equation~\eqref{eq:kappa} is written on the firm's own books. Writing the same transaction
on the counterparty's books exchanges value given for consideration received.

\begin{proposition}[Antisymmetry]
\label{prop:kappa-antisym}
If two parties $i$ and $j$ share the valuation map $v$, and $\kappa_j^{(i)}$ denotes the
credit position towards $j$ on the books of $i$, then
\begin{equation}
\kappa_j^{(i)} = -\,\kappa_i^{(j)}
\label{eq:kappa-antisym}
\end{equation}
\end{proposition}

\begin{proof}
Value given by $i$ to $j$ is value received by $j$ from $i$, and consideration received by
$i$ from $j$ is consideration given by $j$ to $i$. Hence $D_j^{(i)}=P_i^{(j)}$ and
$P_j^{(i)}=D_i^{(j)}$; substituting into~\eqref{eq:kappa} gives the result.
\end{proof}

\begin{remark}[A shared $v$ is required]
Equation~\eqref{eq:kappa-antisym} presumes both parties use the same $v$. As
Remark~\ref{rem:subjectivity} notes, $v$ may be subjective; but the settlement amount $\pi$
is $\mathcal{G}$-measurable and agreed, so any disagreement is confined to the $D$ side.
Differences in acceptance criteria and disputed receivables are the cases in point.
\end{remark}

Antisymmetry has a consequence for aggregation.

\begin{corollary}[Inter-firm credit vanishes on aggregation]
\label{cor:kappa-nets-out}
Let $\mathcal{N}$ be the set of parties in the corporate sector. The sum of credit positions
closed within $\mathcal{N}$ is zero, so
\begin{equation}
\sum_{i\in\mathcal{N}} W^{(i)}
= \kappa_{\text{vs.\ non-corporate}} + \sum_{i\in\mathcal{N}}\mathrm{Inv}^{(i)}
\label{eq:aggregate-W}
\end{equation}
where $\kappa_{\text{vs.\ non-corporate}}$ is the net position of the corporate sector
against households, government and the rest of the world.
\end{corollary}

\begin{proof}
For each pair $(i,j)$ inside $\mathcal{N}$, \eqref{eq:kappa-antisym} gives
$\kappa_j^{(i)}+\kappa_i^{(j)}=0$. Summing the $W$ of Definition~\ref{def:wc} over all
parties in $\mathcal{N}$ cancels the internal pairs and leaves the pairs with the outside
together with inventory.
\end{proof}

\begin{corollary}[Not every stratum can be a net provider of credit]
\label{cor:no-uniform-direction}
Partition the corporate sector into disjoint strata. It is impossible for every stratum to
be a net provider of inter-firm credit. If one stratum is a net provider, another is a net
recipient.
\end{corollary}

\begin{proof}
By Corollary~\ref{cor:kappa-nets-out} the sum of inter-firm positions across strata is zero.
If the net position of every stratum were strictly positive, the sum would be strictly
positive, contradicting zero.
\end{proof}

Corollary~\ref{cor:no-uniform-direction} places a structural constraint on comparisons
across size strata. When the direction of credit is measured by size in
Part~\ref{part:results}, part of the observed sign pattern follows automatically from this
constraint.

\begin{remark}[What the aggregate $\CCC$ measures]
Dividing~\eqref{eq:aggregate-W} by total sales gives the aggregate $\CCC$. Since inter-firm
positions have dropped out of the numerator, \textbf{the aggregate $\CCC$ is not an
indicator of inter-firm credit}. What remains is the position against the non-corporate
sector, plus inventory. Chapter~\ref{part:industry} reports the $\CCC$ of all industries,
and interpreting its level requires this distinction.
\end{remark}

\subsection{Additivity over parallel $\Phi$}
\label{sec:kappa-additive}

One party may run several $\Phi$ in parallel.

\begin{proposition}[Additivity of $\kappa$]
\label{prop:kappa-additive}
If $\Phi^1,\dots,\Phi^K$ are run in parallel on disjoint deliveries
$\delta^1,\dots,\delta^K$, then
\begin{equation}
\kappa = \sum_{k=1}^{K}\kappa^{k},
\qquad
W = \sum_{k=1}^{K} W^{k}
\label{eq:kappa-additive}
\end{equation}
\end{proposition}

\begin{proof}
$D_i$ and $P_i$ in~\eqref{eq:cum} are integrals and therefore additive in the integrand.
Equation~\eqref{eq:kappa} is their difference and preserves additivity.
\end{proof}

What is additive is $\kappa$ and $W$, not $\CCC$. Since $\CCC=W/r$ is a ratio, and ratios
are not additive (Corollary~\ref{cor:ccc-weighted}).

\begin{remark}[The decomposition is not unique]
\label{rem:decomposition-not-unique}
Equation~\eqref{eq:kappa-additive} asserts that a decomposition exists, not that it is
unique. Several families $\{\Phi^k\}$ give the same $\kappa$. The types of
Chapter~\ref{ch:catalog} are a generating set for $\Phi$, not a basis.
\end{remark}

\subsection{Expressing the position as a lag}
\label{sec:tau}

$\kappa_i$ is a difference of cumulants, and what it contains is \textbf{the time gap
between delivery and settlement}. Make this explicit.

\begin{definition}[Settlement lag]
For party $i$, write $\tau_i$ for the average lag from delivery to settlement. Its sign
follows that of $\kappa_i$: settlement later than delivery is positive.
\begin{center}
\begin{tabularx}{\textwidth}{@{}lXl@{}}
\toprule
Sign & State & Examples \\
\midrule
$\tau_i>0$ & settlement later than delivery (the firm extends credit) & receivables, selling on account \\
$\tau_i<0$ & settlement earlier than delivery ($i$ extends credit) & advances received, payables \\
\bottomrule
\end{tabularx}
\captionof{table}{The sign of the settlement lag $\tau_i$ and the corresponding state.}
\label{tab:tau-sign}
\end{center}
\end{definition}

Express the $\kappa_i$ of~\eqref{eq:kappa} through the lag $\tau_i$.

\begin{proposition}[Decomposition of the credit position]
\label{prop:kappa-tau}
In a steady state with a constant rate of delivery, as a time average
\begin{equation}
\bar{\kappa}_i = r\,\tau_i
\label{eq:kappa-tau}
\end{equation}
\end{proposition}

\begin{proof}
If settlement lags delivery uniformly by $\tau_i$, the unsettled cumulant at time $t$ is the
delivery of the most recent $\tau_i$ of time:
\[
\kappa_i(t) = \int_{t-\tau_i}^{t} r(s)\,ds .
\]
If $r$ is constant the right-hand side equals $r\tau_i$. If the lag has a distribution, take
$\tau_i$ to be its mean.
\end{proof}

\begin{remark}[It holds only on average]
\label{rem:kappa-tau-average}
Actual settlement is discrete. Under monthly billing paid at the end of the following month,
$\pi$ is a pulse once a month and $\kappa_i(t)$ moves in a sawtooth.
Equation~\eqref{eq:kappa-tau} holds only for its time average.

In statistics based on period-end balances, if fiscal year-ends cluster in particular months
the phases of the sawtooth line up and the figures can deviate systematically from the mean.
No correction for this bias is made here.
\end{remark}

\begin{example}[$\kappa$ of an annually billed subscription]
A service priced at 1{,}000 yen a month is paid annually, 12{,}000 yen received on 1 January.
Delivery is spread evenly over twelve months, so after $t$ months
\[
\kappa_C(t) = 1{,}000\,t - 12{,}000 .
\]
This is $-12{,}000$ at $t=0$, $-6{,}000$ at $t=6$, and $0$ at $t=12$.
With 100 customers, the mid-period average is $\kappa_C = -600{,}000$ yen.
That amount is interest-free funding.
\end{example}

\section{Decomposing the cash balance}

Use the $\kappa_i$ of~\eqref{eq:kappa} to decompose the cash balance.

\begin{proposition}[Credit decomposition of cash]
\label{prop:cash}
With equity $E(t)$, inventory $\mathrm{Inv}(t)$ and fixed assets $A(t)$,
\begin{equation}
M(t) = E(t) + \sum_{i\in N}\max\{-\kappa_i(t),0\} - \sum_{i\in N}\max\{\kappa_i(t),0\}
- \mathrm{Inv}(t) - A(t)
\label{eq:cash}
\end{equation}
\end{proposition}

\begin{proof}
By the balance-sheet identity, total assets equal total liabilities plus net assets.
Split assets into cash $M$, positive credit positions (receivables, advances paid),
inventory $\mathrm{Inv}$ and fixed assets $A$; map liabilities to negative credit positions
(payables, advances received) and net assets to $E$. Then
\[
M + \sum_i \max\{\kappa_i,0\} + \mathrm{Inv} + A
= \sum_i \max\{-\kappa_i,0\} + E .
\]
Solving for $M$ gives~\eqref{eq:cash}.
\end{proof}

Equation~\eqref{eq:cash} is only a rearrangement of an identity, but it fixes how to read it.
The second term is the total credit drawn from others and the third the total credit extended
to others.

\begin{corollary}[What cash is]
A firm's cash is the total credit drawn from all counterparties, less the credit extended
and the assets fixed in place.
\end{corollary}

\subsection{The net form}
\label{sec:cash-net}

The second and third terms of~\eqref{eq:cash} are written separately in order to show the
gross amounts. Looking only at the net, for any real $x$ we have
$\max\{-x,0\}-\max\{x,0\}=-x$, so the two terms cancel.

\begin{proposition}[Cash is the reverse side of working capital]
\label{prop:cash-net}
With the working capital $W$ of Definition~\ref{def:wc}, \eqref{eq:cash} can be written
\begin{equation}
M(t) = E(t) - W(t) - A(t)
\label{eq:cash-net}
\end{equation}
\end{proposition}

\begin{proof}
Summing $\max\{-\kappa_i,0\}-\max\{\kappa_i,0\}=-\kappa_i$ over $i$ gives
$\sum_i\max\{-\kappa_i,0\}-\sum_i\max\{\kappa_i,0\}=-\sum_i\kappa_i$.
Substitute into~\eqref{eq:cash} and collect $\mathrm{Inv}$ into $W$.
\end{proof}

Checking against Example~\ref{ex:three-sources}: $\sum_i\kappa_i=-600-50+200=-450$ and
$\mathrm{Inv}=0$, so $W=-450$, and $100+450-30=520$ agrees.

Equation~\eqref{eq:cash-net} says the same thing as~\eqref{eq:cash}, but
\textbf{it connects, through $W$, to the $\CCC$ of Chapter~\ref{sec:growth}}.
How the constraint $M\geq 0$ of Section~\ref{sec:objective} bears on the choice of $\Phi$
is settled by way of this form.

\begin{remark}[The identity is a standard one]
Equation~\eqref{eq:cash-net} is the identity used in financial analysis, cash equals equity
minus net operating assets. Its significance here is that $W$ is a function of $\Phi$
through~\eqref{eq:kappa}, so the cash constraint becomes a constraint on $\Phi$.
\end{remark}

\begin{example}[Funding from three sources]
\label{ex:three-sources}
A business started with 1{,}000{,}000 yen of equity receives 6{,}000{,}000 yen of annual
prepayments from customers ($\kappa_C=-6{,}000{,}000$), owes 500{,}000 yen to a cloud
provider ($\kappa_S=-500{,}000$), and has 2{,}000{,}000 yen of receivables from a large
customer ($\kappa_{C'}=+2{,}000{,}000$). There is no inventory and fixed assets are
300{,}000 yen. In units of 10{,}000 yen,
\[
M = 100 + (600+50) - 200 - 0 - 30 = 520 .
\]
Against 1{,}000{,}000 yen of equity, cash on hand is 5{,}200{,}000 yen. The difference of
4{,}200{,}000 yen is credit drawn from customers and suppliers, and it has a due date.
\end{example}

\subsection{Equity and capital providers}
\label{sec:equity}

The $E$ of Proposition~\ref{prop:cash} appears as a residual; its composition is made
explicit here. Equation~\eqref{eq:cash} uses this $E$.

\begin{definition}[Equity]
\label{def:equity}
With paid-in capital $P_K(t)$ and cumulative dividends $\mathrm{Div}(t)$, define
\begin{equation}
E(t) = P_K(t) + \int_0^t \phi(s)\,ds - \mathrm{Div}(t)
\label{eq:equity}
\end{equation}
The second term is retained earnings.
\end{definition}

Definition~\ref{def:equity} links the $\phi$ of Chapter~\ref{sec:surplus} to the $E$ of this
chapter. Substituting the three-way decomposition of $\phi$,
\[
\text{retained earnings} = \int (\phi_{\mathrm{prod}} + \phi_{\mathrm{barg}} + \phi_{\mathrm{cog}})\,ds - \mathrm{Div}
\]
Chapter~\ref{sec:surplus} treats the erosion of $\phi_{\mathrm{cog}}$ by learning and
regulation, but
\textbf{what is eroded is the future flow; what has already accumulated remains in $E$}.

\begin{definition}[Credit position with capital providers]
\label{def:kappa-K}
For a capital provider $i\in K$ with probability measure $\Prob_i$, let
$V_i(t) = \E_i[\text{present value of future cash flows}]$ be the value of the residual
claim, and define
\begin{equation}
\kappa_K^{i}(t) = V_i(t) - P_K(t)
\label{eq:kappa-K}
\end{equation}
\end{definition}

\begin{proposition}[$\kappa_K$ does not enter~\eqref{eq:cash}]
\label{prop:kappa-K-excluded}
Because $\kappa_K^{i}$ depends on the capital provider $i$, it cannot be included in the
identity~\eqref{eq:cash}.
\end{proposition}

\begin{proof}
Equation~\eqref{eq:cash} is a rearrangement of the balance-sheet identity and consists only
of quantities on whose value all parties agree. Receivables and payables agree in amount
between the parties, but the $V_i$ of~\eqref{eq:kappa-K} depends on $\Prob_i$ and does not.
Including a party-dependent quantity would make the identity differ by party, and it would
cease to be an identity.
\end{proof}

\begin{remark}[Three degrees of subjectivity]
\label{rem:subjectivity}
Several unobservable quantities appear in this text, and they differ in character. Using the
$\mathcal{G}$ of Chapter~\ref{sec:info} they fall into three degrees.

\begin{center}
\begin{tabularx}{\textwidth}{@{}llX@{}}
\toprule
Degree & Condition & Examples \\
\midrule
Shared & $\mathcal{G}$-measurable & the settlement amount $\pi$, receivables \\
Private information & $\mathcal{F}_t$-measurable but not $\mathcal{G}$-measurable & the valuation map $v$, effort $a$ \\
Measure-dependent & $\mathcal{F}_t$-measurable but dependent on $\Prob_j$ & the value of the residual claim $V_j$ \\
\bottomrule
\end{tabularx}
\captionof{table}{Three degrees of subjectivity among unobservable quantities.}
\label{tab:subjectivity-levels}
\end{center}

The subjectivity of $v$ is informational asymmetry and can be resolved by disclosure.
The subjectivity of $V_j$ is \textbf{disagreement about the probability measure}; the same
information yields different forecasts, so disclosure does not resolve it.

Two consequences follow. First, the value of $\kappa_K^{i}$ does not agree between the
parties and so differs in character from the other $\kappa_i$. Second,
\textbf{it is precisely that disagreement that makes equity investment possible}: an
investor invests because they value $V$ more highly than the manager does. This has the same
structure as the heterogeneity of discount factors in Proposition~\ref{prop:discount}.
\end{remark}

The maturity of capital is discussed in Remark~\ref{rem:tau-K}.

\begin{remark}[$\tau_K$ and insolvency]
\label{rem:tau-K}
For borrowing, $\tau_K$ is fixed as the repayment date; for equity it is not fixed. Capital
is the longest-dated credit and has no maturity.

Making dividends obligatory would fix $\kappa_K$, but~\eqref{eq:constraint} would then break
immediately upon non-payment. Leaving them non-obligatory as a residual claim, and forcing
liquidation at the insolvency time of Section~\ref{sec:objective}, avoids this.
Bankruptcy is the mechanism that fixes $\kappa_K$.
\end{remark}

\begin{remark}[Funds held on behalf of others]
Among positions with $\kappa_i<0$, funds that legally belong to someone else --- escrow, for
instance --- occupy the same place in~\eqref{eq:cash} but differ in character. They drain
rapidly when volume falls, and must not be treated as equivalent to other negative $\kappa$.
\end{remark}

\chapter{Information Structure and the Objective: When a Contract Can Be Written, and the Path Constraint}
\label{sec:info}

\section{Verifiable shared information}
\label{sec:verifiable}

Write $\mathcal{G}_i\subseteq\mathcal{F}$ for the information party $i$ can observe, and
$\mathcal{G}=\bigcap_i \mathcal{G}_i$ for the shared information verifiable between the
parties.

\begin{proposition}[Implementability]
\label{prop:implementable}
A necessary condition for a settlement schedule $\pi$ to be implementable as a contract is
that $\pi$ be $\mathcal{G}$-measurable.
\end{proposition}

\begin{proof}
Suppose $\pi$ is not $\mathcal{G}$-measurable. Then there exist two states $\omega,\omega'$
indistinguishable under $\mathcal{G}$ with $\pi(\omega)\neq\pi(\omega')$. The payment
obligation depends on which state obtains, yet the parties cannot jointly verify which one
does. They therefore cannot reach agreement on the amount, and no third party can enforce
performance.
\end{proof}

This condition constrains the design of $\Phi$ severely. When the effort level $a$ is not in
$\mathcal{G}$, no contract conditioned on $a$ can be written and only the observable outcome
$y$ can be used.

\begin{center}
\begin{tabularx}{\textwidth}{@{}llX@{}}
\toprule
What is observable & Contract form & Examples \\
\midrule
Effort $a$ & input-linked & hourly wages, cost plus \\
Only the outcome $y$ & outcome-contingent & success fees, royalties, deductibles \\
Both, partially & mixed & fixed plus commission \\
\bottomrule
\end{tabularx}
\captionof{table}{What is observable and the corresponding contract form.}
\label{tab:observability-contracts}
\end{center}

This is a standard result of contract theory (\cite{holmstrom1979}) and is not a
contribution of this text.

\begin{remark}[What $\mathcal{G}$ bounds is degree of freedom (4)]
\label{rem:G-limits-dof4}
Proposition~\ref{prop:implementable} bears directly on degree of freedom (4), state
dependence, in Section~\ref{sec:dof}. $\pi$ may be made to depend on $\omega$ only within
$\mathcal{G}$; distinctions among $\omega$ that lie outside $\mathcal{G}$ cannot be written
into a contract.

\textbf{So $\mathcal{G}$ bounds the codomain of $\Phi$, not its domain.} Any delivery
$\delta$ may be chosen; what is constrained is how $\pi$ may be made to move with it.

The constraint appears in the assignment order of Part~\ref{part:catalog}. The family
decided by degree of freedom (4) is family 5 (outcome- and state-contingent), and that family
cannot be implemented unless the $y$ in $\pi=h(y)$ is $\mathcal{G}$-measurable. The
impossibility result for family 7 in Chapter~\ref{ch:protocol}
(Proposition~\ref{prop:beneficiary}) has the same shape.
\end{remark}

\begin{remark}[Endogeneity of $\mathcal{G}$]
\label{rem:G-endogenous}
This chapter treats $\mathcal{G}$ as given. In practice \textbf{an operator can invest in
$\mathcal{G}$}: follow-up surveys make the outcome $y$ observable and widen the range of
verifiable information.

The staffing-service disclosures treated in Section~\ref{sec:jinzai-structure} include a
field for the number of placements whose separation could not be ascertained, which shows
that the breadth of $\mathcal{G}$ differs by operator and is itself a choice.

Proposition~\ref{prop:implementable} should therefore be read as a short-run condition with
$\mathcal{G}$ held fixed. In the long run $\mathcal{G}$ becomes part of $\Phi$.
\end{remark}

\begin{remark}[A distinction of direction]
Risk transfer comes in two kinds that run in opposite directions. Success fees and
royalties leave risk with the agent \emph{because effort is unobservable}. Cost plus returns
risk to the principal \emph{because effort is observable but the outcome is uncertain}.
Grouping both under ``risk sharing'' erases the difference.
\end{remark}

\begin{example}[Optimal strength of outcome contingency]
\label{ex:bstar}
Consider the linear contract $\pi = \alpha + b\,y$. Writing $\gamma$ for the agent's
absolute risk aversion and $\sigma^2$ for the noise variance of the outcome
($r$ is reserved for the rate of sales in Chapter~\ref{sec:growth}, hence $\gamma$ here),
the optimal outcome-contingency coefficient takes the form
\begin{equation}
b^{\star} = \frac{1}{1+\gamma\,\sigma^{2}\,k}
\label{eq:bstar}
\end{equation}
where $k$ is a constant relating to the marginal cost of effort (\cite{holmstrom1991}).
The larger $\sigma^2$, the smaller $b^\star$: \textbf{the noisier the outcome, the weaker
the outcome contingency should be}.

As a numerical illustration, with $\gamma k = 1$: $b^\star=0.67$ at $\sigma^2=0.5$,
$b^\star=0.33$ at $\sigma^2=2$, and $b^\star=0.10$ at $\sigma^2=9$.
\end{example}

\section{The objective: a cash constraint imposed path by path}
\label{sec:objective}

The $M(t,\omega)$ in the constraint is the cash balance of Proposition~\ref{prop:cash}.

\begin{definition}[The optimization problem]
With period profit $\phi_t$ and discount factor $\beta$,
\begin{align}
\max_{\Phi}\quad & \E\!\left[\sum_{t\in T}\beta^{t}\phi_t\right] \label{eq:obj}\\
\text{s.t.}\quad & M(t,\omega)\geq 0 \qquad \forall t\in T,\ \forall\omega\in\Omega \label{eq:constraint}
\end{align}
\end{definition}

The character of~\eqref{eq:constraint} is the centre of this text.

\begin{proposition}[The path constraint does not commute with expectation]
\label{prop:pathwise}
$\E[M(t)]\geq 0$ does not imply~\eqref{eq:constraint}. Moreover the insolvency time
$\Tins=\inf\{t: M(t)<0\}$ is absorbing, and no $\phi_t$ for $t>\Tins$ is realized.
Because $\tau_i$ denotes the settlement lag, the insolvency time is written $\Tins$.
\end{proposition}

\begin{proof}
A counterexample suffices for the first claim. If $M(t)$ takes $+100$ with probability
$1/2$ and $-100$ with probability $1/2$, then $\E[M(t)]=0\geq 0$, yet the constraint is
violated on a path of probability $1/2$.

The second claim holds because insolvency occurs where $M(\Tins)<0$ and no further delivery
$\delta$ can be performed. The objective should therefore be written
$\E\bigl[\sum_{t<\Tins}\beta^t\phi_t\bigr]$, and $\Tins$ depends on $\Phi$.
\end{proof}

\begin{corollary}[Insolvency while profitable]
\label{cor:insolvency-with-profit}
Even if $\phi_t>0$ for every $t$, \eqref{eq:constraint} can be violated. Insolvency while
profitable is thus a failure of the constraint~\eqref{eq:constraint}, not of the
objective~\eqref{eq:obj}.
\end{corollary}

\begin{proof}
By~\eqref{eq:cash-net} in Chapter~\ref{sec:growth}, $M=E-W-A$. A positive $\phi$ raises $E$
through~\eqref{eq:equity}, but if the increase in $W$ exceeds it, $M$ falls.
Proposition~\ref{prop:depletion} gives the depletion time in that case.
\end{proof}

The condition for insolvency while profitable is rewritten as $g>g^\star$ in
Chapter~\ref{sec:growth} (Remark~\ref{rem:level-and-change}).

\begin{example}[Insolvency while profitable, numerically]
\label{ex:insolvency}
A business with margin $m=10\%$ grows from 100 million yen of annual revenue to 200 million.
Its profit is 20 million yen.

If credit to customers runs at 20\% of revenue and inventory at 5\%, the sum of the credit
positions of~\eqref{eq:kappa} and inventory is 25 million yen at 100 million of revenue and
50 million at 200 million. The increase of 25 million yen leaves in cash.

Against a profit of 20 million yen, 5 million is short. The books show a profit and the cash
is short, and~\eqref{eq:constraint} can be violated.
\end{example}

\section{Heterogeneity of discount factors}

\begin{proposition}[Gain from rearranging the settlement schedule]
\label{prop:discount}
If $\beta_C<\beta_0$ (the customer discounts the future more heavily than the firm), there
exists a rearrangement of the timing of $\pi$ alone, leaving the delivery $\delta$ unchanged,
that improves both parties' subjective valuations simultaneously.
\end{proposition}

\begin{proof}
Let the current contract call for a payment $p$ at time $0$, and consider deferring it to a
payment $p'$ at time $1$. The customer's valuation changes from $-p$ to $-\beta_C p'$, so
the customer improves if $\beta_C p' < p$, that is $p' < p/\beta_C$. The firm's valuation
changes from $p$ to $\beta_0 p'$, so the firm improves if $\beta_0 p' > p$, that is
$p' > p/\beta_0$. Since $\beta_C<\beta_0$ we have $p/\beta_0 < p/\beta_C$, so the interval
is non-empty. Any $p'$ in it improves both.
\end{proof}

\begin{example}[The gain from instalments]
Suppose the customer's annualized discount rate is 20\% ($\beta_C=0.833$) and the firm's cost
of funds is 5\% ($\beta_0=0.952$). For an item with a cash price of 100{,}000 yen, a price
$p'$ payable in one year improves both parties for $105{,}000 < p' < 120{,}000$. Taking
$p'=112{,}000$: the firm receives $112{,}000\times 0.952 = 106{,}624$ yen in present value,
a gain of 6{,}624 yen; the customer values it subjectively at
$112{,}000\times 0.833=93{,}296$ yen, a gain of 6{,}704 yen.
\end{example}

\begin{remark}[Two readings]
$\beta_C<\beta_0$ can have two different causes. For a customer facing a liquidity
constraint, present funds are objectively more valuable. Alternatively the cause may be
cognitive, as with hyperbolic discounting (\cite{laibson1997}). Meeting the first and
exploiting the second are almost indistinguishable from outside, yet their durability
(Chapter~\ref{sec:surplus}) is opposite. This framework cannot settle the distinction.
\end{remark}

\chapter{The Growth Constraint: \texorpdfstring{$\CCC$}{CCC} and the Self-Financeable Growth Rate}
\label{sec:growth}

\section{The cash conversion cycle}
\label{sec:ccc-def}

Let $r(t)$ be the rate of sales and $W(t)$ the working capital of Definition~\ref{def:wc}.
Since $W$ is a sum of the $\kappa_i$ of~\eqref{eq:kappa}, it depends on $\Phi$.

\begin{definition}[Cash conversion cycle]
In a steady state, define
\begin{equation}
\CCC = \frac{W}{r}
\label{eq:ccc}
\end{equation}
Its dimension is time.
\end{definition}

\section{The fundamental inequality}
\label{sec:fundamental}

Substituting~\eqref{eq:ccc} into~\eqref{eq:cash-net} of Proposition~\ref{prop:cash-net},
the cash balance can be written
\begin{equation}
M(t) = E(t) - \CCC\cdot r(t) - A(t)
\label{eq:M-ccc}
\end{equation}
The constraint $M\geq 0$ of Section~\ref{sec:objective} therefore takes the following form.

\begin{proposition}[The fundamental inequality]
\label{prop:fundamental}
The cash constraint~\eqref{eq:constraint} is equivalent to
\begin{equation}
\CCC\cdot r(t,\omega) \ \leq\ E(t) - A(t)
\label{eq:fundamental}
\end{equation}
\end{proposition}

\begin{proof}
Substitute $W=\CCC\cdot r$ into~\eqref{eq:cash-net} and rearrange $M\geq 0$.
\end{proof}

Equation~\eqref{eq:fundamental} reads: \textbf{working capital cannot exceed equity net of
fixed assets}. The left-hand side is fixed by $\Phi$, the right-hand side by the capital
structure. The rest of this chapter, and each chapter of Part~\ref{part:catalog}, treats a
special case of this inequality.

\begin{corollary}[Upper bound on scale]
\label{cor:rmax}
When $\CCC>0$, the feasible rate of sales is bounded above by
\begin{equation}
r \ \leq\ r^{\max} = \frac{E-A}{\CCC}
\label{eq:rmax}
\end{equation}
When $\CCC\leq 0$ and $E\geq A$,
\eqref{eq:fundamental} holds for every $r$.
\end{corollary}

The dimension is yen per year, matching the rate of sales.
Where $g^\star$ is a bound on the growth rate, $r^{\max}$ is \textbf{a bound on the level}.
The two are separate constraints and neither implies the other
(Remark~\ref{rem:level-and-change}).

\begin{corollary}[Without capital, only $\Phi$ with $\CCC\leq 0$ can be chosen]
\label{cor:zero-equity}
When $E=0$, \eqref{eq:fundamental} requires $\CCC\cdot r\leq -A\leq 0$, that is
$\CCC\leq 0$.
\end{corollary}

Corollary~\ref{cor:zero-equity} is the starting point of Chapter~\ref{ch:solo}.
What remains available to a party without capital is only those $\Phi$ in which settlement
precedes delivery.

\section{Growth and free cash flow}
\label{sec:fcf}

The fundamental inequality is a constraint on the level. Change is treated next.

\begin{proposition}[Growth and free cash flow]
\label{prop:fcf}
Writing $m$ for the margin and $I$ for capital expenditure, and
\textbf{taking $\CCC$ to be constant and independent of $m$},
\begin{equation}
\FCF = m\,r - \CCC\cdot\dot{r} - I
\label{eq:fcf}
\end{equation}
\end{proposition}

\begin{proof}
Free cash flow is gross margin less the increase in working capital and capital expenditure:
\[
\FCF = m\,r - \dot{W} - I .
\]
By~\eqref{eq:ccc}, $W=\CCC\cdot r$, and if $\CCC$ is constant then
$\dot W = \CCC\cdot\dot r$. Substituting gives~\eqref{eq:fcf}.
\end{proof}

\begin{remark}[$\FCF$ is the time derivative of the cash balance]
\label{rem:fcf-is-mdot}
Differentiate~\eqref{eq:M-ccc}. By~\eqref{eq:equity}, $\dot E=\phi-\dot{\mathrm{Div}}$, and
fixed assets satisfy $\dot A = I - d\,r$ (investment less depreciation). With no dividends
and no external funding,
\[
\dot M = \phi + d\,r - \CCC\,\dot r - I = \FCF
\]
So $\FCF$ is not an independently introduced quantity but \textbf{the time derivative
of~\eqref{eq:cash-net}}. The $g^\star$ of the next section is therefore the boundary of
$\dot M\geq 0$.
\end{remark}

Solve~\eqref{eq:fcf} for $g$. The $\CCC$ of~\eqref{eq:ccc} appears in the denominator.

\begin{corollary}[Self-financeable growth rate]
\label{cor:gstar}
Writing $d$ for depreciation as a fraction of sales and $g=\dot r/r$ for the growth rate,
the condition $\FCF\geq 0$ is
% The only boxed equation in this text. \boxed is fine for PDF, but tex4ht emits
% <menclose notation="box">, which MathML Core lacks: Chrome and Safari draw no box and
% Firefox clips its top edge against the fraction. In HTML an empty span marks the spot and
% CSS draws the box on the adjacent equation table. A bare <div> inserted mid-paragraph
% breaks the <p> nesting, make4ht's domfilter abandons XML parsing, the whole file is
% re-serialized, and < inside formulas is corrupted. A balanced empty element avoids this.
\ifdefined\HCode
\HCode{<span class="keyeq-mark"></span>}%
\begin{equation}
g \ \leq\ g^{\star} = \frac{m + d - I/r}{\CCC}
\label{eq:gstar}
\end{equation}
\else
\begin{equation}
\boxed{\ g \ \leq\ g^{\star} = \frac{m + d - I/r}{\CCC}\ }
\label{eq:gstar}
\end{equation}
\fi
When $\CCC<0$ there is no $g^\star$: $\FCF>0$ at any growth rate.
\end{corollary}

\begin{proof}
If $m$ is the operating margin, depreciation has already been deducted, yet depreciation
involves no cash outlay. The cash-basis gross margin is therefore $(m+d)\,r$.
Correcting~\eqref{eq:fcf} gives $\FCF = (m+d)\,r - \CCC\,\dot r - I$.
Substituting $\dot r = g\,r$ and dividing by $r$,
$\FCF/r = (m+d) - \CCC\,g - I/r$.
Solving $\FCF\geq 0$ for $g$ gives~\eqref{eq:gstar}.
When $\CCC<0$ the left-hand side is positive for every $g>0$.
\end{proof}

\section{Level and change are separate constraints}
\label{sec:level-and-change}

\begin{remark}[The constraint on the level and the constraint on change are independent]
\label{rem:level-and-change}
Corollary~\ref{cor:rmax} follows from $M\geq 0$ and Corollary~\ref{cor:gstar} from
$\dot M\geq 0$. Neither implies the other. Even with $M(0)>0$, if $g>g^\star$ then $M$
falls and~\eqref{eq:constraint} is violated in finite time.
\end{remark}

\begin{proposition}[Time to cash depletion]
\label{prop:depletion}
Let $I=0$, let $g$ be constant, and let $r(t)=r_0e^{gt}$. If
$g>g^\star=(m+d)/\CCC$, then starting from $M(0)=M_0>0$ the cash runs out at
\begin{equation}
\Tins = \frac{1}{g}\,
\ln\!\left(1+\frac{g\,M_0}{(\CCC\,g-m-d)\,r_0}\right)
\label{eq:depletion}
\end{equation}
and $\Tins\to\infty$ as $g\downarrow g^\star$.
\end{proposition}

\begin{proof}
By Remark~\ref{rem:fcf-is-mdot}, $\dot M = a\,r_0e^{gt}$ with $a=(m+d)-\CCC g$.
Integrating, $M(t)=M_0+(a r_0/g)(e^{gt}-1)$. Since $g>g^\star$ implies $a<0$, $M$ falls, and
solving $M(\Tins)=0$ gives~\eqref{eq:depletion}. As $g\to g^\star$, $a\to 0$ and the
argument of the logarithm diverges.
\end{proof}

Equation~\eqref{eq:depletion} puts a value on the insolvency time that
Proposition~\ref{prop:pathwise} introduced only as a symbol, and fixes when the insolvency
of Corollary~\ref{cor:insolvency-with-profit} actually arrives.

\begin{example}[$g^\star$ is a threshold, not a cliff]
\label{ex:depletion}
With $m=5\%$, $d=2\%$, $\CCC=0.25$ years, $M_0=100$ and $r_0=1{,}000$, we have
$g^\star=28\%$.
\begin{center}
\begin{tabular}{@{}lr@{}}
\toprule
Growth rate & $\Tins$ \\
\midrule
$g^\star+0.001$ & 16.8 years \\
$g^\star+0.01$ & 8.7 years \\
$g^\star+0.05$ & 3.9 years \\
$g^\star+0.2$ & 1.4 years \\
$g^\star+0.5$ & 0.6 years \\
\bottomrule
\end{tabular}
\captionof{table}{Excess over $g^\star$ and the time to cash depletion.}
\label{tab:depletion}
\end{center}
A slight excess leaves more than a decade; a large one exhausts the cash within a year.
\textbf{What matters is not whether $g^\star$ has been exceeded but by how much.}
\end{example}

\begin{remark}[$\CCC$ is also an amplifier of demand variation]
\label{rem:ccc-amplifier}
Equation~\eqref{eq:M-ccc} is linear in $r$ with coefficient $-\CCC$. If $E$ and $A$ are
independent of $r$ in the short run,
\[
\operatorname{Var}[M] = \CCC^{2}\operatorname{Var}[r]
\]
So \textbf{$\CCC$ not only fixes the bound on the level; it amplifies variation in demand
into variation in cash.} The larger a sector's $\CCC$, the more readily
\eqref{eq:constraint} is violated for the same variation in demand.

Because no process is specified for $r$, no $\Prob(\Tins\leq t)$ is given here. That the
amplification coefficient is $\CCC$ can be stated without specifying the process.
\end{remark}

\begin{remark}[No assumption that $I \approx d\,r$]
\label{rem:no-steady-investment}
With a constant asset base, $I = d\,r$ and~\eqref{eq:gstar} reduces to $m/\CCC$. That
assumption is not made here. The measurements of Section~\ref{sec:gstar-measured} put $I/r$
between $-1.8\%$ and $6.7\%$ across industries and size classes, which does not agree
with $d$.
\end{remark}

\begin{remark}[Divergence when $\CCC$ is small]
\label{rem:gstar-divergence}
Equation~\eqref{eq:gstar} diverges as $\CCC\to 0$. Even for positive $\CCC$, if it is only a
few days then a small movement in the numerator moves $g^\star$ a great deal.

At $\CCC=5$ days, for instance, a one-point change in the numerator moves $g^\star$ by about
73 points.
\textbf{In sectors where $\CCC$ is small, $g^\star$ does not work as an indicator.}
Section~\ref{sec:gstar-measured} shows this by reporting standard deviations alongside.
\end{remark}

$g^\star$ is the rate of growth that can be financed internally and is in substance
identical to the self-financeable growth rate of \cite{churchill2001}.
The relation between $\FCF$ and the growth rate is shown in Figure~\ref{fig:growth}.

\begin{figure}[htbp]
\centering
\insertfig{growth-fcf}
\caption{Growth rate and free cash flow. The vertical intercept is the numerator of~\eqref{eq:gstar}; the sign of $\CCC$ fixes the sign of the slope.}
\label{fig:growth}
\end{figure}

\begin{example}[The magnitudes $g^\star$ actually takes]
Consider a business with operating margin $m=5\%$, depreciation $d=2.5\%$ of sales, and
capital expenditure $I/r=3.5\%$ --- close to the all-industry figures measured in
Section~\ref{sec:gstar-measured}. The numerator is
\[
m + d - I/r = 5.0 + 2.5 - 3.5 = 4.0\%
\]
Computing $g^\star$ for several values of $\CCC$:

\begin{center}
\begin{tabular}{@{}rrrr@{}}
\toprule
$\CCC$ (days) & $\CCC$ (years) & $g^\star$ & For reference: $m/\CCC$ \\
\midrule
120 & 0.329 & 12.2\% & 15.2\% \\
90 & 0.247 & 16.2\% & 20.3\% \\
60 & 0.164 & 24.3\% & 30.4\% \\
30 & 0.082 & 48.7\% & 60.9\% \\
\textminus 30 & \textminus 0.082 & does not exist & does not exist \\
\bottomrule
\end{tabular}
\captionof{table}{$g^\star$ and $m/\CCC$ for several values of $\CCC$ (illustrative).}
\label{tab:gstar-example}
\end{center}

A business with $\CCC=120$ days necessarily needs outside funding to grow faster than 12\% a
year. A business with $\CCC=-30$ days accumulates cash the faster it grows.
\textbf{The sign fixes the qualitative constitution; the magnitude fixes the quantitative
room.}

The right-hand column gives the value when $d$ and $I$ are ignored. Here $d<I/r$, so
$m/\CCC$ overstates $g^\star$. \textbf{For businesses investing above depreciation, an
assessment using $m$ alone is optimistic.}
\end{example}

\begin{remark}[Domain of definition]
Equation~\eqref{eq:ccc} presumes a steady state and is undefined immediately after founding,
where $r\approx 0$. This is taken up again in Section~\ref{sec:age}.
\end{remark}

\section{\texorpdfstring{$\CCC$}{CCC} is a sum of lags}
\label{sec:ccc-as-tau}

Substituting Proposition~\ref{prop:kappa-tau} into~\eqref{eq:ccc} eliminates $r$.

\begin{proposition}[Structure of $\CCC$]
\label{prop:ccc-structure}
In a steady state,
\begin{equation}
\CCC = \sum_i \tau_i + \tau_{\mathrm{inv}},
\qquad \tau_{\mathrm{inv}} \equiv \frac{\mathrm{Inv}}{r}
\label{eq:ccc-tau}
\end{equation}
That is, $\CCC$ does not depend on the flow; it is a sum of lags.
\end{proposition}

\begin{proof}
From $W=\sum_i\kappa_i+\mathrm{Inv}$ and~\eqref{eq:kappa-tau},
$W = r\sum_i\tau_i + \mathrm{Inv}$. Divide both sides by $r$.
Since~\eqref{eq:kappa-tau} holds as a time average
(Remark~\ref{rem:kappa-tau-average}), this proposition is likewise a statement about
averages.
\end{proof}

Three consequences follow from~\eqref{eq:ccc-tau}. All are deductions and use no
observation.

\begin{corollary}[Source of differences in level]
\label{cor:ccc-level}
$\tau_i$ is the quantity specified by degree of freedom (1) of Section~\ref{sec:dof}, the
order of $\operatorname{supp}(\pi)$ and $\operatorname{supp}(\delta)$. The level of $\CCC$
is therefore fixed by $\Phi$.
\end{corollary}

\begin{corollary}[Cross-sectional comparison is meaningless]
\label{cor:no-cross-comparison}
Comparing the level of $\CCC$ across objects with different $\Phi$ is comparing different
contractual forms, and carries no information about demand or efficiency.
\end{corollary}

\begin{corollary}[Aggregation is an $r$-weighted mean]
\label{cor:ccc-weighted}
In the parallel arrangement of Proposition~\ref{prop:kappa-additive}, writing $r_k$ for the
rate of sales of each part,
\begin{equation}
\CCC = \sum_k w_k\,\CCC_k, \qquad w_k = \frac{r_k}{\sum_j r_j}
\label{eq:ccc-weighted}
\end{equation}
so $\CCC$ is not additive.
\end{corollary}

\begin{proof}
From $W=\sum_k W_k$ (Proposition~\ref{prop:kappa-additive}) and $W_k=\CCC_k r_k$,
$\CCC=\sum_k \CCC_k r_k/\sum_j r_j$.
\end{proof}

Equation~\eqref{eq:ccc-weighted} means that computing $\CCC$ for a firm running $\Phi$ from
several families in parallel yields a weighted mean of terms with opposite signs. This is
the ground for taking $\Phi$, not the firm, as the unit of observation.

\begin{corollary}[Decomposition of change]
\label{cor:ccc-change}
Differencing~\eqref{eq:ccc-tau} splits it as
\begin{equation}
\Delta\CCC = \underbrace{\sum_i \Delta\tau_i}_{\text{change in contracts}}
 + \underbrace{\Delta\tau_{\mathrm{inv}}}_{\text{change in operations}}
\label{eq:ccc-change}
\end{equation}
In~\eqref{eq:ccc-change}, $\Delta\tau_i\neq 0$ means a change in contractual terms, whereas
$\Delta\tau_{\mathrm{inv}}$ arises independently of contracts.
\end{corollary}

\section{The path constraint and the peak}
\label{sec:peak}

Equation~\eqref{eq:constraint} is imposed for every $t$ and every $\omega$. Yet, as
Remark~\ref{rem:kappa-tau-average} notes, $\bar\kappa_i=r\tau_i$ is a time average.
\textbf{What binds is the peak, not the average.}

\begin{proposition}[Working capital at the peak]
\label{prop:peak}
With a constant rate of sales, a billing period $T_b$ and a payment lag $\tau_i$, the time
average of $\kappa_i(t)$ is $r(\tau_i+T_b/2)$ and its maximum is $r(\tau_i+T_b)$.
Equation~\eqref{eq:fundamental} therefore becomes
\begin{equation}
r \ \leq\ \frac{E-A}{\CCC + T_b/2}
\label{eq:rmax-peak}
\end{equation}
\end{proposition}

\begin{proof}
At each billing date $kT_b$ the deliveries of $[(k-1)T_b,kT_b)$ are invoiced and settled at
$kT_b+\tau_i$. The unsettled cumulant is the sum of what has accrued since the last invoice
and what has been invoiced but not paid; it is a sawtooth with minimum $r\tau_i$ just after a
billing date and maximum $r(\tau_i+T_b)$ just before one. Its average over a cycle is
$r(\tau_i+T_b/2)$. The peak exceeds the $\CCC$ measured as an average by $T_b/2$, so
evaluating the left-hand side of~\eqref{eq:fundamental} at the peak
gives~\eqref{eq:rmax-peak}.
\end{proof}

With monthly billing, $T_b/2$ is about fifteen days.
\textbf{Estimating the capital required from a $\CCC$ computed as an average understates it
by half a month.} The measurements in Part~\ref{part:results} rest on period-end balances or
within-period averages and do not include this difference.

\section{Comparative statics}
\label{sec:comparative-statics}

Every quantity in this chapter is a function of parameters. Setting the derivatives side by
side shows the designer of $\Phi$ what to move and by how much.

\begin{center}
\begin{tabularx}{\textwidth}{@{}llX@{}}
\toprule
Quantity & Derivative & Content \\
\midrule
$g^\star$ & $\partial\ln g^\star/\partial\tau_i = -1/\CCC$ & independent of $i$ \\
$g^\star$ & $\partial g^\star/\partial m = 1/\CCC$ & improving the margin bites harder the smaller $\CCC$ is \\
Eq.~\eqref{eq:rmax} & $\partial\ln r^{\max}/\partial\tau_i = -1/\CCC$ & the same coefficient as for $g^\star$ \\
Eq.~\eqref{eq:rmax} & $\partial\ln r^{\max}/\partial E = 1/(E-A)$ & raising equity bites harder the larger $A$ is \\
$\operatorname{Var}[M]$ & $\partial\operatorname{Var}[M]/\partial\CCC = 2\CCC\operatorname{Var}[r]$ & Remark~\ref{rem:ccc-amplifier} \\
\bottomrule
\end{tabularx}
\captionof{table}{Comparative statics of the quantities in this chapter.}
\label{tab:comparative-statics}
\end{center}

\begin{proposition}[Lags are perfect substitutes at the margin]
\label{prop:tau-substitution}
By~\eqref{eq:ccc-tau}, $\partial\ln g^\star/\partial\tau_i = -1/\CCC$ for every $i$,
independently of $i$.
\end{proposition}

\begin{proof}
From $g^\star=(m+d-I/r)/\CCC$ and $\CCC=\sum_j\tau_j+\tau_{\mathrm{inv}}$ we have
$\partial\CCC/\partial\tau_i=1$, hence
$\partial\ln g^\star/\partial\tau_i = -\partial\ln\CCC/\partial\tau_i = -1/\CCC$.
\end{proof}

\textbf{Collecting one day earlier and paying one day later are equivalent for $g^\star$.}
By the third row of Table~\ref{tab:comparative-statics}, the same coefficient appears for
$r^{\max}$. Whether one looks at the level or at change, the design guidance for $\Phi$ is
the same.

Corollary~\ref{cor:no-cross-comparison} supplies a deductive ground for the limitation
stated in Remark~\ref{rem:level-vs-change}. That the quantities here are defined as levels
is a presentational choice, but \textbf{the impossibility of cross-sectional comparison is a
consequence of the definitions}.

\begin{remark}[No comparable quantity can be constructed]
\label{rem:no-normalization}
Corollary~\ref{cor:no-cross-comparison} states why comparison fails; it does not supply a
way to compare.

The natural candidate is the deviation from the $\tau_i^{\mathrm{contract}}$ that the
contract specifies,
\[
\tau_i^{\mathrm{observed}} - \tau_i^{\mathrm{contract}}
\]
This quantity does not depend on $\Phi$, so it can be compared across industries, and it
measures whether performance follows the contract.

But $\tau_i^{\mathrm{contract}}$ is the payment term of an individual contract and does not
exist in published statistics. \textbf{The theoretical problem is solved and the measurement
problem takes its place} --- category (iv), access constraint, in the taxonomy of
Section~\ref{sec:obstacle-taxonomy}.
\end{remark}

\begin{remark}[Independence of $m$ and $\CCC$]
\label{rem:m-ccc-independence}
Corollary~\ref{cor:gstar} gives $g^\star$ as a ratio but does not guarantee that numerator
and denominator move independently.

The measurements of Section~\ref{sec:gstar-measured} show $\CCC$ rising 49\% and $m$ rising
91\% between fiscal 2000 and fiscal 2024 --- \textbf{both increasing}. If holding more
working capital and earning a higher margin are positively related, a change in $g^\star$
cannot be read as a loosening of the constraint.

This text does not formalize the relation between the two. $g^\star$ is reported as the
value of a ratio and given no causal interpretation.
\end{remark}

\chapter{Surplus and Layers: Source, Durability, Mutability}
\label{sec:surplus}

\section{The three-way decomposition of surplus}
\label{sec:surplus-decomp}

The surplus $\phi$ announced in Remark~\ref{rem:Phi-vs-phi} is defined here. Period profit
is decomposed by source into three terms; the decomposition is an organization contributed
by this text. The third term is the cognitive surplus of~\eqref{eq:cog} and can be positive
when degree of freedom (5) of~\eqref{eq:phi} is present. The first two are defined in
Definition~\ref{def:phi-prod}.
\begin{equation}
\phi = \phi_{\mathrm{prod}} + \phi_{\mathrm{barg}} + \phi_{\mathrm{cog}}
\label{eq:surplus}
\end{equation}

\begin{center}
\begin{tabularx}{\textwidth}{@{}llX@{}}
\toprule
Symbol & Source & Character and decay \\
\midrule
$\phi_{\mathrm{prod}}$ & advantage in marginal cost $c(\delta)$ & technical advantage; does not reduce anyone else's share; eroded by competition but reproducible through investment \\
$\phi_{\mathrm{barg}}$ & the map $\Phi$ & transfer of shares; zero-sum with $\sum_i \phi_i = 0$; constrained by the institutional layer \\
$\phi_{\mathrm{cog}}$ & divergence between $\dbar$ and $\delta$ & eroded by learning, regulation and better interfaces, and in principle not recovered \\
\bottomrule
\end{tabularx}
\captionof{table}{The three-way decomposition of surplus $\phi$, with the source and durability of each component.}
\label{tab:surplus-decomposition}
\end{center}

\begin{remark}[Where the return to bearing risk belongs]
\label{rem:phi-risk}
A party that makes a business of the transfer and pooling of Section~\ref{sec:three-ops} ---
an insurer or a payment processor, say --- earns a return for bearing risk. Where does it
belong in~\eqref{eq:surplus}?

The part by which variance is actually reduced through pooling is $\phi_{\mathrm{prod}}$.
By Proposition~\ref{prop:pooling}, raising $n$ removes the $\sigma^2/n$ term; this is a
technical advantage and reduces nobody else's share.

The part corresponding to the limit $\rho\sigma^2$ of~\eqref{eq:pooling}, by contrast,
cannot be removed. The return for bearing it is a transfer and satisfies
$\sum_i\phi_i=0$; it belongs to $\phi_{\mathrm{barg}}$.

In practice the two are not separated and appear as a single rate. Where convenient this
text writes $\phi_{\mathrm{risk}}$ (Section~\ref{sec:pred-J}), which is not a new component
but the undivided sum of $\phi_{\mathrm{prod}}$ and $\phi_{\mathrm{barg}}$.
\end{remark}

\begin{definition}[Production surplus and bargaining surplus]
\label{def:phi-prod}
Using the marginal cost $c(\delta)$ of Definition~\ref{def:cost} and the price $p^{\ast}$
that would obtain without market power, define
\begin{align}
\phi_{\mathrm{prod}} &= \bigl(p^{\ast} - c(\delta)\bigr)\,\delta \\
\phi_{\mathrm{barg}} &= \bigl(p - p^{\ast}\bigr)\,\delta
\end{align}
where $p$ is the actual price. The character of $p^{\ast}$ is discussed in
Remark~\ref{rem:counterfactual-price}.
\end{definition}

\begin{remark}[$p^{\ast}$ is a counterfactual]
\label{rem:counterfactual-price}
Definition~\ref{def:phi-prod} requires $p^{\ast}$, which is not observed. It is a
counterfactual price posited for a state without market power, and has the same structure as
the standard procedure for measuring market power in industrial organization.

Separating $\phi_{\mathrm{prod}}$ from $\phi_{\mathrm{barg}}$ therefore requires an
additional assumption identifying $p^{\ast}$.

One assumption suffices: \textbf{that a control is observable which is identical except for
market power}. Given such a control, its price can be read as $p^{\ast}$ and the difference
attributed to $\phi_{\mathrm{barg}}$. Chapter~\ref{sec:platform-fee} satisfies this in the
form of two rates on the same platform that differ only in function. Where the assumption
fails, this text goes no further than the sum
$\phi_{\mathrm{prod}}+\phi_{\mathrm{barg}}$.
\end{remark}

Model the decay of cognitive surplus as a function of the transaction frequency $\nu$:
\begin{equation}
\dot{\phi}_{\mathrm{cog}} = -\lambda(\nu)\,\phi_{\mathrm{cog}},
\qquad \lambda'(\nu)>0
\label{eq:decay}
\end{equation}

\begin{proposition}[Implication of the decay]
The solution of~\eqref{eq:decay} is
$\phi_{\mathrm{cog}}(t)=\phi_{\mathrm{cog}}(0)e^{-\lambda(\nu)t}$, with half-life
$\ln 2/\lambda(\nu)$. The higher the transaction frequency, the faster cognitive surplus
disappears.
\end{proposition}

\begin{proof}
Direct integration of a linear first-order ordinary differential equation. Solving
$\phi(t)=\phi(0)e^{-\lambda t}=\phi(0)/2$ gives $t=\ln 2/\lambda$. Since
$\lambda'(\nu)>0$, the half-life is decreasing in $\nu$.
\end{proof}

\begin{example}[Frequency and half-life]
Suppose $\lambda$ is proportional to the number of transactions a year, $\lambda = 0.02\nu$.
\begin{center}
\begin{tabular}{@{}llrr@{}}
\toprule
Good & Transactions a year $\nu$ & $\lambda$ & Half-life \\
\midrule
Daily necessities & 200 & 4.0 & 0.17 years \\
Monthly service & 12 & 0.24 & 2.9 years \\
Annual contract (insurance, memberships) & 1 & 0.02 & 34.7 years \\
\bottomrule
\end{tabular}
\captionof{table}{Transaction frequency $\nu$ and the half-life of cognitive surplus (illustrative).}
\label{tab:decay-halflife-example}
\end{center}

\section{Surplus and cash}
\label{sec:surplus-cash}

$\phi$ affects cash only through~\eqref{eq:equity}, and the three components do not affect
it in the same way.

\begin{proposition}[Recognizing cognitive surplus does not move cash]
\label{prop:cog-cash-neutral}
Suppose $p\,\dbar$ is received under a prepaid contract and delivery ends at
$\delta<\dbar$. When the residual obligation lapses at expiry and $\phi_{\mathrm{cog}}$ is
recognized, the $M$ of~\eqref{eq:cash-net} does not change.
\end{proposition}

\begin{proof}
With the valuation $v(\delta)=p\,\delta$, the credit position before recognition is
$\kappa_C = p\delta - p\dbar = -p(\dbar-\delta) = -\phi_{\mathrm{cog}}$.
On recognition $\kappa_C$ becomes $0$, so $W$ rises by $\phi_{\mathrm{cog}}$.
At the same time $E$ rises by the same amount by~\eqref{eq:equity}.
In~\eqref{eq:cash-net} the two cancel.
\end{proof}

\textbf{Cognitive surplus generates no cash at the moment of recognition.}
The cash arrived when the prepayment was received, and its effect there appeared not as
$\phi_{\mathrm{cog}}$ but as $\kappa_C<0$. The $\phi_{\mathrm{cog}}$ of~\eqref{eq:cog} and
the undelivered balance of Chapter~\ref{sec:kappa} are two cross-sections of the same
transaction.

\begin{center}
\begin{tabularx}{\textwidth}{@{}lX@{}}
\toprule
Component & When the cash moves \\
\midrule
$\phi_{\mathrm{prod}}$, $\phi_{\mathrm{barg}}$ & at settlement, lagged by $\tau_C$ \\
$\phi_{\mathrm{cog}}$ & at contracting, not at recognition \\
\bottomrule
\end{tabularx}
\captionof{table}{The three components of~\eqref{eq:surplus} and when the cash moves.}
\label{tab:surplus-cash-timing}
\end{center}

\begin{remark}[This agrees with the accounting treatment]
Proposition~\ref{prop:cog-cash-neutral} says the same thing as the transfer from contract
liabilities (breakage) under revenue-recognition standards. Its significance here is the
translation into the language of $\kappa$: the transfer moves both $E$ and $W$
in~\eqref{eq:cash-net} by the same amount, so they cancel.
\end{remark}

\begin{corollary}[The equity that can be built from cognitive surplus is bounded]
\label{cor:cog-bounded}
Under the decay~\eqref{eq:decay}, the total equity that can be accumulated from
$\phi_{\mathrm{cog}}$ is bounded by
\begin{equation}
\int_0^{\infty}\phi_{\mathrm{cog}}(0)\,e^{-\lambda(\nu)s}\,ds
= \frac{\phi_{\mathrm{cog}}(0)}{\lambda(\nu)}
\label{eq:cog-total}
\end{equation}
Since $\lambda'(\nu)>0$, the higher the transaction frequency the lower the bound.
\end{corollary}

Substituting~\eqref{eq:cog-total} into Corollary~\ref{cor:rmax}, for a business with no
source of surplus other than cognitive surplus,
\begin{equation}
r \ \leq\ \frac{E(0) + \phi_{\mathrm{cog}}(0)/\lambda(\nu) - A}{\CCC}
\label{eq:rmax-cog}
\end{equation}
\textbf{Equation~\eqref{eq:rmax-cog} shows that the transaction frequency fixes the upper
bound on the scale a business can reach.} Equation~\eqref{eq:decay} does not stop at a
discussion of half-lives; it connects to the constraint of Chapter~\ref{sec:growth}.
The coefficient values have no empirical basis and only illustrate the structure of the
ratio. What follows is an ordering: goods bought daily retain almost no cognitive surplus,
while contracts on an annual cycle retain it for a long time.
\end{example}

\section{Layers: five levels ordered by mutability}
\label{sec:layers}

The foregoing is organized into five layers ordered by mutability. The higher the layer, the
more readily it is rewritten by outside forces.

\begin{center}
\begin{tabularx}{\textwidth}{@{}llX@{}}
\toprule
Layer & Object & Mutability \\
\midrule
Institutional & regulation, commercial custom, interface norms & changes exogenously and rewrites the four layers below at once \\
Cognitive & $\dbar-\delta$ & eroded monotonically by learning and regulation \\
Credit & the allocation of $\kappa_i$ & depends on bargaining power; constrained by the institutional layer \\
Timing & the relation between $\operatorname{supp}(\pi)$ and $\operatorname{supp}(\delta)$ & largely determined by the physical layer \\
Physical & $C_{\mathrm{prod}}$, marginal cost, asset intensity & technology and physical law; given in the short run \\
\bottomrule
\end{tabularx}
\captionof{table}{Five layers ordered by mutability.}
\label{tab:five-layers}
\end{center}

This order coincides with the order of durability of surplus. A $\phi$ that rests on the
cognitive and institutional layers is eroded continuously from two directions, legal change
and customer learning. An advantage originating in the physical layer is harder to erode.

\begin{remark}[The institutional layer does not act uniformly]
\label{rem:institution-modes}
``The institutional layer imposes constraints'' covers at least three distinct forms.

\begin{center}
\begin{tabularx}{\textwidth}{@{}lXl@{}}
\toprule
Form & Effect & Example in this text \\
\midrule
Entry requirement & removes the feasibility of a $\Phi$ & deposit obligations, licensing (Section~\ref{sec:bootstrap}) \\
Price regulation & hollows out an option & capped fees (Part~\ref{part:results}) \\
Manipulation of competitive conditions & erodes $\phi_{\mathrm{barg}}$ indirectly & litigation pressure, fee deregulation (Remark~\ref{rem:institution-shift}) \\
\bottomrule
\end{tabularx}
\captionof{table}{Three forms in which the institutional layer imposes constraints.}
\label{tab:institution-modes}
\end{center}

The three erode differently. The first two make a particular $\Phi$ unavailable; the third
leaves $\Phi$ in place and lowers the price.
\textbf{The dominant form in recent years is the third; forms that regulate price directly
are in decline} (Section~\ref{sec:institution-mode}).

Saying that something is ``eroded by the institutional layer'' therefore says nothing about
how fast or how surely. Chapter~\ref{sec:platform-fee} treats a case in which fees arising
from market power can be separated from fees arising from function, but what is observed
there is a single reduction of the third form, not continuous decay.
\end{remark}
